% =====================================================================
%  INTRODUCCIÓN · cómo usar la guía, la cocina de la web, el laboratorio
% =====================================================================
\clearpage
\renewcommand{\seccionactual}{CÓMO USAR ESTA GUÍA}
\section*{Cómo usar esta guía}
\addcontentsline{toc}{section}{Cómo usar esta guía}

Hoy vas a aprender \textbf{PHP desde cero}, aunque nunca hayas escrito una línea de código de servidor. La clase está
armada como un juego: \textbf{7 niveles}, cada uno con su desafío y sus puntos de experiencia (XP).

\subsection*{Cada nivel tiene siempre las mismas partes}
\renewcommand{\arraystretch}{1.1}
\begin{tabularx}{\linewidth}{>{\leavevmode\sffamily\bfseries\color{fuegooscuro}}p{4.6cm} X}
  La idea en simple & Una comparación con la vida real. Si solo leés una cosa, que sea esta. \\
  Míralo funcionar & Abrís la página en el navegador y ves el resultado. \\
  El código explicado & El archivo, \textbf{línea por línea}, en castellano. \\
  Modificalo & Cambiás cosas pequeñas y ves qué pasa. Así se aprende de verdad. \\
  El reto & Escribís algo vos. La solución está en la carpeta \texttt{soluciones}, pero \textbf{intentá antes}. \\
  Errores típicos & Lo que le pasa a todo el mundo, y cómo arreglarlo. \\
  Comprobá que entendiste & Preguntas cortas. Las respuestas están en el Apéndice C. \\
\end{tabularx}

\subsection*{El método: ver $\rightarrow$ modificar $\rightarrow$ crear}
\begin{center}
\begin{tikzpicture}[font=\sffamily\small]
  \node[draw=cianjp,fill=cianjp!8,rounded corners=2mm,line width=1pt,minimum width=4.4cm,minimum height=1.5cm,align=center] (a) at (0,0) {\textbf{1 · VER}\\el ejemplo funcionando};
  \node[draw=orojp,fill=orojp!10,rounded corners=2mm,line width=1pt,minimum width=4.4cm,minimum height=1.5cm,align=center] (b) at (5.4,0) {\textbf{2 · MODIFICAR}\\cambiar una línea y ver qué pasa};
  \node[draw=rojojp,fill=rojojp!8,rounded corners=2mm,line width=1pt,minimum width=4.4cm,minimum height=1.5cm,align=center] (c) at (10.8,0) {\textbf{3 · CREAR}\\escribir algo propio};
  \draw[-{Stealth},thick,grisjp] (a) -- (b); \draw[-{Stealth},thick,grisjp] (b) -- (c);
\end{tikzpicture}
\end{center}

\subsection*{Tu rango de cocinero}
\renewcommand{\arraystretch}{1.05}
\begin{tabularx}{\linewidth}{l X >{\bfseries}l}
\toprule
\textbf{Nivel} & \textbf{Qué aprendés} & XP \\
\midrule
0 · El gancho & Por qué existe PHP y dónde se ejecuta & 50 \\
1 · Hablar & \texttt{echo}, variables, comillas & 100 \\
2 · Decidir y repetir & \texttt{if}, arrays, \texttt{foreach} & 150 \\
3 · Escuchar & Formularios con \texttt{\$\_GET} y \texttt{\$\_POST} & 150 \\
4 · Recordar & PHP + MySQL con PDO & 200 \\
5 · El museo & Los 5 errores que todos cometen & 100 \\
X · Jefe final & PHP en tu proyecto FORJA WEB & 250 \\
\midrule
 & \textbf{Total} & 1000 \\
\bottomrule
\end{tabularx}
\vspace{1mm}

\begin{center}
\begin{tikzpicture}[font=\sffamily\small]
  \foreach \i/\c/\t/\r in {0/grisjp/Ayudante de cocina/0--299,1/orojp/Cocinero/300--599,2/fuego/Chef/600--899,3/rojojp/Chef ejecutivo/900--1000}{
    \node[draw=\c,fill=\c!10,rounded corners=2mm,line width=1pt,minimum width=3.7cm,minimum height=1.3cm,align=center] at ({\i*4.05},0) {\textbf{\t}\\\texttt{\r} XP};}
\end{tikzpicture}
\end{center}

% =====================================================================
\clearpage
\renewcommand{\seccionactual}{LA COCINA DE LA WEB}
\section*{La cocina de la web}
\addcontentsline{toc}{section}{La cocina de la web}

\subsection*{HTML muestra. PHP piensa.}
Hasta ahora hiciste páginas con \textbf{HTML}: dicen \textbf{siempre lo mismo}. Si escribís «Son las 10:00», dice
«Son las 10:00» aunque sean las tres de la tarde. HTML sabe \emph{mostrar}, pero no sabe \emph{pensar}.

\textbf{PHP} es un lenguaje de programación que corre \textbf{en el servidor}. Antes de mandarte la página, PHP puede
calcular la hora, decidir qué saludo poner, leer datos de una base de datos, recibir lo que escribiste en un
formulario\ldots\ y recién entonces \textbf{arma el HTML} y te lo envía.

\begin{analogia}
Pensá en un restaurante. Vos (el \textbf{navegador}) le pedís un plato al \textbf{mozo} (\textbf{Apache}, el servidor web).
Si pedís algo ya hecho, como una gaseosa (un archivo \texttt{.html}), el mozo te lo trae tal cual. Si pedís algo que hay
que cocinar (un archivo \texttt{.php}), el mozo se lo pasa al \textbf{cocinero} (\textbf{PHP}), que busca ingredientes en la
\textbf{despensa} (\textbf{MySQL}), cocina y le entrega el plato terminado al mozo. \textbf{Vos nunca entrás a la cocina}:
solo ves el plato servido (el HTML).
\end{analogia}

\begin{center}
\begin{tikzpicture}[font=\sffamily\small,
  c/.style={draw=#1,fill=#1!8,rounded corners=2mm,line width=1pt,minimum width=3.1cm,minimum height=1.45cm,align=center}]
  \node[c=grisjp] (n) at (0,0) {\textbf{Navegador}\\(el cliente)};
  \node[c=cianjp] (a) at (4.4,0) {\textbf{Apache}\\(el mozo)};
  \node[c=fuego] (p) at (8.8,0) {\textbf{PHP}\\(el cocinero)};
  \node[c=lilaneon!80!black] (m) at (13.2,0) {\textbf{MySQL}\\(la despensa)};
  \draw[-{Stealth},thick,grisjp] (n.north east) to[bend left=18] node[above,font=\scriptsize]{pide \texttt{hora.php}} (a.north west);
  \draw[-{Stealth},thick,cianjp] (a.north east) to[bend left=18] node[above,font=\scriptsize]{«cociná esto»} (p.north west);
  \draw[-{Stealth},thick,lilaneon!80!black] (p.north east) to[bend left=18] node[above,font=\scriptsize]{pide datos} (m.north west);
  \draw[-{Stealth},thick,verdejp] (p.south west) to[bend left=18] node[below,font=\scriptsize]{HTML terminado} (a.south east);
  \draw[-{Stealth},thick,verdejp] (a.south west) to[bend left=18] node[below,font=\scriptsize]{la página} (n.south east);
\end{tikzpicture}
\end{center}

\subsection*{La idea más importante del día}
\begin{marco}[fuego]{PHP se ejecuta en el SERVIDOR}
El navegador \textbf{nunca} ve tu código PHP. Solo recibe el \textbf{resultado}: HTML. Por eso PHP puede guardar
secretos (como la contraseña de la base de datos) que nadie desde afuera puede leer.
\end{marco}

\subsection*{¿Por qué PHP?}
\begin{itemize}
  \item Según W3Techs (agosto de 2026), PHP se usa en alrededor del \textbf{70\,\%} de los sitios web cuyo lenguaje de servidor se conoce: \textbf{7 de cada 10}.
  \item \textbf{WordPress}, con el que se hacen muchísimos blogs, tiendas y sitios de empresas, está hecho en PHP.
  \item Es de los más fáciles para empezar: hoy, en tres horas, vas a guardar datos en una base de datos.
\end{itemize}

% =====================================================================
\clearpage
\renewcommand{\seccionactual}{ENCENDER LA COCINA}
\section*{Encender la cocina (tu laboratorio)}
\addcontentsline{toc}{section}{Encender la cocina (tu laboratorio)}

Todo corre en \textbf{Docker}, así que no tenés que instalar XAMPP ni configurar nada.
El laboratorio tiene \textbf{tres contenedores}:

\renewcommand{\arraystretch}{1.3}
\begin{tabularx}{\linewidth}{>{\ttfamily\bfseries}l l X}
\toprule
\normalfont\bfseries Contenedor & \textbf{Qué es} & \textbf{Dónde lo ves} \\
\midrule
fogon-web & Apache + PHP 8.3 (mozo y cocinero) & \texttt{http://localhost:8090} \\
fogon-db  & MySQL 8.4 (la despensa) & \texttt{127.0.0.1:3308} (Workbench) \\
fogon-pma & phpMyAdmin (la ventana a la despensa) & \texttt{http://localhost:8091} \\
\bottomrule
\end{tabularx}

\parte{Paso 1 · Requisitos}
\begin{itemize}
  \item \textbf{Docker Desktop} abierto (la ballena tiene que estar quieta, no animada).
  \item \textbf{VS Code} (o cualquier editor de texto; el Bloc de notas también sirve).
  \item La carpeta \texttt{operacion-fogon} que comparte el profe, copiada en tu Escritorio.
\end{itemize}

\parte{Paso 2 · Encender}
Abrí la carpeta \texttt{operacion-fogon}, hacé \textbf{clic derecho en un espacio vacío $\rightarrow$ Abrir en Terminal} y escribí:
\begin{terminal}[PowerShell]{la primera vez tarda unos minutos (descarga imágenes)}
Set-ExecutionPolicy -Scope Process Bypass
.\scripts\iniciar.ps1
\end{terminal}
El script levanta todo, espera a que MySQL esté listo y abre el navegador en \texttt{http://localhost:8090}.
Tenés que ver el \textbf{mapa de la cocina} con el cocinero y la despensa \textbf{encendidos}.

\captura[0.85\linewidth]{imagenes/index.png}

\parte{Paso 3 · Abrir tus recetas en VS Code}
En VS Code: \textbf{Archivo $\rightarrow$ Abrir carpeta} $\rightarrow$ elegí la carpeta \texttt{operacion-fogon}. Tus archivos
están dentro de \texttt{www}:
\begin{terminal}[estructura]{carpeta operacion-fogon}
www/
  00-gancho/       nivel 0: HTML vs PHP
  01-hablar/       nivel 1: echo y variables
  02-decidir/      nivel 2: if, arrays, foreach
  03-escuchar/     nivel 3: formularios GET y POST
  04-recordar/     nivel 4: PHP + MySQL
  05-errores/      el museo de los 5 errores
  06-reto-final/   jefe final: formulario de contacto
  soluciones/      soluciones de los retos
db/init.sql        la despensa: tablas y datos iniciales
docker-compose.yml el plano de la cocina
\end{terminal}

\parte{Paso 4 · El ciclo que vas a repetir 100 veces}
\begin{center}
\begin{tikzpicture}[font=\sffamily\small,
  c/.style={draw=#1,fill=#1!8,rounded corners=2mm,line width=1pt,minimum width=3.6cm,minimum height=1.3cm,align=center}]
  \node[c=cianjp] (a) at (0,0) {\textbf{1 · Editás}\\en VS Code};
  \node[c=orojp] (b) at (4.8,0) {\textbf{2 · Guardás}\\\texttt{Ctrl + S}};
  \node[c=fuego] (c) at (9.6,0) {\textbf{3 · Recargás}\\\texttt{F5} en el navegador};
  \draw[-{Stealth},thick,grisjp] (a) -- (b); \draw[-{Stealth},thick,grisjp] (b) -- (c);
  \draw[-{Stealth},thick,verdejp] (c.south) to[bend left=22] node[below,font=\scriptsize]{y otra vez} (a.south);
\end{tikzpicture}
\end{center}
No hace falta reiniciar nada: la carpeta \texttt{www} está «conectada» al contenedor. Guardás y el servidor ya tiene el cambio.

\begin{ojo}[Si algo no anda]
\textbf{El puerto está ocupado} (8090, 8091 o 3308): en \texttt{docker-compose.yml} cambiá el número de la \emph{izquierda}, por ejemplo \texttt{"8095:80"}.
\textbf{«La despensa está cerrada»}: MySQL tarda unos segundos la primera vez; esperá y recargá.
\textbf{Docker no funciona en tu PC}: usá el plan B, \texttt{.\textbackslash scripts\textbackslash plan-b-sin-base.ps1}, que levanta solo PHP y alcanza para los niveles 0 a 3 y el museo.
\end{ojo}
